\clearpage
\phantomsection
\addcontentsline{toc}{chapter}{Simulado e Revisão Integrada do Módulo 5}
\chapter*{Simulado e Revisão Integrada: Módulo 5}
\markboth{SIMULADO DO MÓDULO 5}{SIMULADO DO MÓDULO 5}

\noindent\tagPratica\ \tagEmenta
\vspace{6pt}

\begin{boxConceito}[Instruções do Simulado Integrado (Capítulos 13, 14 e 15)]
Este simulado integra os conteúdos das Semanas 14 a 16 (Bioengenharia de Água e Solo, Legislação Ambiental/Código Florestal/Licenciamento/PRAD e Manual Prático do Diagnóstico Ambiental Aplicado), consolidando o encerramento da disciplina CHF662.
\begin{itemize}[leftmargin=1.5em, itemsep=2pt]
  \item \textbf{Tempo sugerido:} $90$ minutos sem consulta ao gabarito.
  \item \textbf{Estrutura:} $8$ questões objetivas autorais no estilo ENADE e concursos ambientais (valor: $0{,}75$ ponto cada $= 6{,}0$ pontos) e $2$ questões dissertativas aplicadas (valor: $2{,}0$ pontos cada $= 4{,}0$ pontos).
  \item Ao concluir, confira o \textbf{Gabarito Comentado}, o \textbf{Espelho Dissertativo} e preencha a \textbf{Tabela de Diagnóstico de Revisão} na página seguinte.
\end{itemize}
\end{boxConceito}

\section*{Parte I: Questões Objetivas (Estilo ENADE e Concursos Ambientais)}

\begin{enumerate}[leftmargin=1.8em, itemsep=8pt]
  \item \textbf{[Cap. 13, Bioengenharia e Critério de Mohr--Coulomb]} Na estabilização de taludes e margens fluviais por meio da bioengenharia de água e solo (Gray; Sotir, 1996; Durlo; Sutili, 2005), as raízes das plantas contribuem diretamente para o aumento da resistência ao cisalhamento do solo segundo o critério de Mohr--Coulomb ($\tau = c' + c_r + \sigma'\tan\phi'$). A parcela $c_r$ e um efeito hidrológico benéfico típico da cobertura vegetal correspondem, respectivamente, a:
  \begin{alternativas}
    \item Pressão neutra positiva dos poros e aumento da velocidade do escoamento superficial.
    \item Coesão adicional conferida pela malha radicular e interceptação da chuva associada à redução da poropressão via evapotranspiração.
    \item Peso próprio do dossel arbóreo sobre a crista do talude e selamento impermeável do horizonte A.
    \item Redução do ângulo de atrito interno ($\phi'$) e impermeabilização completa da zona vadosa.
    \item Substituição da tensão efetiva ($\sigma'$) por escoamento laminar concentrado.
  \end{alternativas}

  \item \textbf{[Cap. 13, Dimensionamento de Paliçadas Permeáveis]} Para conter o aprofundamento de uma ravina com extensão longitudinal $L = 45{,}0\text{ m}$, largura média de seção $b = 1{,}5\text{ m}$ e declividade longitudinal do leito $S = 8\%$ ($0{,}08\text{ m/m}$), projetaram-se paliçadas permeáveis de madeira e bambu com altura útil do vertedouro $h = 0{,}60\text{ m}$. Adotando o espaçamento horizontal geométrico $E = h / S$ e aproximando cada cunha de sedimento retida por um prisma triangular ($V_{\text{cunha}} = \tfrac{1}{2}\,b\,h\,E$), o espaçamento $E$, o número de estruturas ($n = L / E$) e o volume total de sedimentos retido ($V_{\text{total}} = n \cdot V_{\text{cunha}}$) são, respectivamente:
  \begin{alternativas}
    \item $E = 5{,}0\text{ m}$, $n = 9$ paliçadas e $V_{\text{total}} = 20{,}25\text{ m}^3$.
    \item $E = 7{,}5\text{ m}$, $n = 6$ paliçadas e $V_{\text{total}} = 13{,}50\text{ m}^3$.
    \item $E = 7{,}5\text{ m}$, $n = 6$ paliçadas e $V_{\text{total}} = 20{,}25\text{ m}^3$.
    \item $E = 12{,}0\text{ m}$, $n = 4$ paliçadas e $V_{\text{total}} = 40{,}50\text{ m}^3$.
    \item $E = 6{,}0\text{ m}$, $n = 7$ paliçadas e $V_{\text{total}} = 18{,}90\text{ m}^3$.
  \end{alternativas}

  \item \textbf{[Cap. 13, Biomantas, Geotêxteis e Espécies Nativas]} Em um talude de corte com solo exposto e elevada declividade no Semiárido Baiano, recomenda-se a aplicação de biomanta biodegradável associada a semeadura de gramíneas/leguminosas e estacas vivas adaptadas. A função imediata da biomanta nas primeiras semanas após a implantação é:
  \begin{alternativas}
    \item Substituir permanentemente o sistema radicular das árvores adultas durante décadas sem necessidade de germinação vegetal.
    \item Absorver a energia cinética do impacto das gotas de chuva (\textit{splash}), reduzir o escoamento superficial, reter umidade/sementes e proteger o solo até o estabelecimento da vegetação.
    \item Impedir totalmente a infiltração da água da chuva no solo, promovendo escoamento hortoniano máximo.
    \item Eliminar o lençol freático profundo por impermeabilização asfáltica da superfície.
    \item Introduzir espécies exóticas invasoras de alta inflamabilidade para acelerar a degradação do solo.
  \end{alternativas}

  \item \textbf{[Cap. 14, PNMA e Licenciamento Ambiental Trifásico]} À luz da Política Nacional do Meio Ambiente (Lei nº 6.938/1981) e da sistemática brasileira de licenciamento ambiental trifásico (Sánchez, 2013), assinale a alternativa \textbf{correta}:
  \begin{alternativas}
    \item A responsabilidade civil pela reparação de danos ambientais no Brasil é subjetiva, exigindo comprovação de dolo do empreendedor para obrigar à recuperação da área degradada.
    \item A Licença de Operação (LO) é emitida antes da Licença Prévia (LP) para agilizar a supressão vegetal de áreas prioritárias.
    \item A Licença Prévia (LP) atesta a viabilidade ambiental da localização e concepção do empreendimento; a Licença de Instalação (LI) autoriza o início das obras conforme projetos aprovados; e a Licença de Operação (LO) autoriza o funcionamento após verificação do cumprimento das condicionantes anteriores.
    \item A Licença de Instalação (LI) dispensa qualquer programa de monitoramento ambiental ou PRAD durante a fase de obras.
    \item O Estudo de Impacto Ambiental (EIA/RIMA) deve ser elaborado somente após a conclusão integral das obras e o encerramento das atividades do empreendimento.
  \end{alternativas}

  \item \textbf{[Cap. 14, Código Florestal: APP e Reserva Legal]} Um imóvel rural de $100{,}0\text{ ha}$ situado no bioma Caatinga (fora da Amazônia Legal) é atravessado em linha reta por um curso d'água natural perene com $8\text{ m}$ de largura regular ao longo de $1.000\text{ m}$ de extensão, com ambas as margens dentro da propriedade. Considerando o Código Florestal (Lei nº 12.651/2012, Arts. 4º e 12) e sem computar a APP no percentual de Reserva Legal (Art. 15), a largura da APP em cada margem, a área total de APP marginal nesse trecho e a área mínima de Reserva Legal (RL) do imóvel são, respectivamente:
  \begin{alternativas}
    \item $15\text{ m}$ em cada margem; $3{,}0\text{ ha}$ de APP; e $20{,}0\text{ ha}$ de Reserva Legal ($20\%$).
    \item $30\text{ m}$ em cada margem; $6{,}0\text{ ha}$ de APP; e $20{,}0\text{ ha}$ de Reserva Legal ($20\%$).
    \item $50\text{ m}$ em cada margem; $10{,}0\text{ ha}$ de APP; e $35{,}0\text{ ha}$ de Reserva Legal ($35\%$).
    \item $30\text{ m}$ em cada margem; $3{,}0\text{ ha}$ de APP; e $80{,}0\text{ ha}$ de Reserva Legal ($80\%$).
    \item $100\text{ m}$ em cada margem; $20{,}0\text{ ha}$ de APP; e $20{,}0\text{ ha}$ de Reserva Legal ($20\%$).
  \end{alternativas}

  \item \textbf{[Cap. 14, Plano de Recuperação de Áreas Degradadas (PRAD)]} Sobre os conceitos e etapas operacionais de um PRAD (Brasil/MMA, 2016; Sánchez, 2013), assinale a alternativa \textbf{INCORRETA}:
  \begin{alternativas}
    \item A \textbf{restauração} busca aproximar o ecossistema degradado de sua composição e estrutura originais, enquanto a \textbf{reabilitação} visa restabelecer funções ecológicas e produtivas compatíveis com um uso sustentável, sem necessariamente retornar ao estado pré-distúrbio idêntico.
    \item O diagnóstico prévio da área degradada (solo, drenagem, erosão e remanescentes vegetais) é condição indispensável antes da definição do modelo de revegetação.
    \item Todo PRAD consistente deve incluir cronograma físico de execução e indicadores de monitoramento pós-implantação (como sobrevivência de mudas, cobertura do solo e estabilização erosiva).
    \item O isolamento da área contra fatores de degradação contínua (como fogo e pisoteio intensivo de gado) compõe a estratégia básica de proteção das mudas implantadas.
    \item A entrega do relatório inicial do PRAD ao órgão ambiental encerra definitivamente a responsabilidade do empreendedor, dispensando acompanhamento de campo ou replantio nos anos seguintes.
  \end{alternativas}

  \item \textbf{[Cap. 15, Bases Abertas e Cartografia do Trabalho Final]} Durante a elaboração dos mapas temáticos do Trabalho Prático Final (Semana 16) no QGIS, uma equipe precisa delimitar uma microbacia, avaliar sua declividade e cruzar com a cobertura vegetal histórica. O conjunto mais adequado de bases abertas e padrão cartográfico é:
  \begin{alternativas}
    \item MDE SRTM/TOPODATA para delimitação e declividade, MapBiomas para séries históricas de uso/cobertura, SICAR/CAR para perímetros rurais e projeção \textbf{SIRGAS 2000 / UTM} com os seis elementos cartográficos obrigatórios.
    \item Apenas croqui desenhado à mão livre sem coordenadas nem escala gráfica.
    \item Coordenadas geográficas em graus decimais sem projeção plana para calcular diretamente áreas de APP em metros quadrados.
    \item Exclusivamente dados censitários demográficos urbanos para estimar a declividade das vertentes.
    \item Mapas sem legenda, sem seta Norte e sem referência à fonte dos dados para economizar espaço na página.
  \end{alternativas}

  \item \textbf{[Cap. 15, Matriz de Indicadores e Cálculo do $IPI$]} Na Matriz de Indicadores do Trabalho Prático Final, um setor de encosta cultivada apresentou Fragilidade Emergente de Ross $FE = 4$, Índice de Qualidade do Solo $IQS = 0{,}52$ e Pressão/Impacto PEIR normalizado $P_{\text{norm}} = 0{,}70$. Adotando a equação $IPI = 0{,}35\,F_{\text{norm}} + 0{,}35\,(1 - IQS) + 0{,}30\,P_{\text{norm}}$, com $F_{\text{norm}} = (FE - 1)/4$, o valor do $IPI$ e sua classe de prioridade são:
  \begin{alternativas}
    \item $0{,}480$ (Prioridade Média).
    \item $0{,}585$ (Prioridade Média).
    \item $0{,}6405$ (Prioridade Alta).
    \item $0{,}7455$ (Prioridade Alta).
    \item $0{,}8250$ (Prioridade Crítica).
  \end{alternativas}
\end{enumerate}

\section*{Parte II: Questões Dissertativas Aplicadas}

\begin{enumerate}[leftmargin=1.8em, itemsep=8pt, start=9]
  \item \textbf{[Integração Caps. 13 e 14, Bioengenharia, Código Florestal e PRAD]} Um empreendimento rodoviário em processo de regularização ambiental causou supressão indevida de $40\text{ m}$ de mata ciliar na margem de um rio de $15\text{ m}$ de largura, desencadeando solapamento do barranco e erosão em sulcos no talude adjacente.
  \begin{enumerate}[label=(\alph*), leftmargin=1.6em, itemsep=2pt]
    \item Enquadre a faixa de APP exigida pelo Código Florestal (Lei nº 12.651/2012) para esse curso d'água e explique por que a obrigação de reparar o dano aplica-se independentemente da existência de culpa ou dolo (Lei nº 6.938/1981).
    \item Proponha um arranjo técnico de \textbf{bioengenharia de água e solo} e \textbf{PRAD} para estabilizar a margem e recompor a APP, citando pelo menos duas técnicas complementares e dois indicadores de monitoramento pós-implantação.
  \end{enumerate}
  \espacoResolucao[4.0cm]

  \item \textbf{[Integração Cap. 15, Diagnóstico Ambiental Aplicado e $IPI$]} Na preparação do Trabalho Prático Final (Semana 16, $5{,}0$ pontos), um grupo levantou os seguintes indicadores para dois setores de uma microbacia:
  \begin{itemize}[leftmargin=1.4em, itemsep=1.5pt]
    \item \textbf{Setor 1 (Cabeceira degradada):} $FE = 5$, $IQS = 0{,}32$ e $P_{\text{norm}} = 0{,}85$.
    \item \textbf{Setor 2 (Terço inferior com manejo parcial):} $FE = 2$, $IQS = 0{,}76$ e $P_{\text{norm}} = 0{,}30$.
  \end{itemize}
  \begin{enumerate}[label=(\alph*), leftmargin=1.6em, itemsep=2pt]
    \item Calcule o $IPI$ de cada setor (usando $w_F = 0{,}35$, $w_S = 0{,}35$, $w_P = 0{,}30$) e classifique ambos quanto à prioridade de intervenção.
    \item Indique três elementos cartográficos obrigatórios que não podem faltar nos mapas temáticos do grupo (em SIRGAS 2000/UTM) e explique qual técnica de intervenção é mais compatível com o Setor 1.
  \end{enumerate}
  \espacoResolucao[4.2cm]
\end{enumerate}

\clearpage
\section*{Gabarito Comentado e Espelho de Correção: Simulado do Módulo 5}

\begin{boxGabarito}[Gabarito Comentado: Parte I (Questões Objetivas 1 a 8)]
\begin{enumerate}[leftmargin=1.6em, itemsep=4pt]
  \item \textbf{(B)}, No critério de Mohr--Coulomb adaptado, $c_r$ representa a coesão adicional conferida pela resistência à tração da malha de raízes; hidrologicamente, a vegetação intercepta parte da chuva e reduz a poropressão da água nos poros por evapotranspiração.
  \item \textbf{(C)}, Espaçamento: $E = h / S = 0{,}60 / 0{,}08 = \mathbf{7{,}5\text{ m}}$. Número de estruturas: $n = 45{,}0 / 7{,}5 = \mathbf{6}$ paliçadas. Volume de cada cunha: $V_{\text{cunha}} = 0{,}5 \times 1{,}5 \times 0{,}60 \times 7{,}5 = 3{,}375\text{ m}^3$. Volume total retido: $V_{\text{total}} = 6 \times 3{,}375 = \mathbf{20{,}25\text{ m}^3}$.
  \item \textbf{(B)}, A biomanta atua como proteção superficial imediata contra o \textit{splash}, reduz a velocidade do escoamento, retém umidade e sementes e estabiliza o leito até que as raízes da vegetação implantada assumam o papel estrutural permanente.
  \item \textbf{(C)}, Apenas (C) descreve corretamente a sequência trifásica do licenciamento ambiental brasileiro: LP (viabilidade locacional e conceitual), LI (autorização de instalação/obra) e LO (autorização de funcionamento após atendimento das condicionantes).
  \item \textbf{(B)}, Pela Lei nº 12.651/2012 (Art. 4º, I), cursos d'água com menos de $10\text{ m}$ de largura exigem faixa marginal de \textbf{$30\text{ m}$} em cada margem. Área de APP nas duas margens: $2 \times 30\text{ m} \times 1.000\text{ m} = 60.000\text{ m}^2 = \mathbf{6{,}0\text{ ha}}$. Fora da Amazônia Legal (Caatinga), a Reserva Legal mínima é de $20\%$, logo $0{,}20 \times 100{,}0 = \mathbf{20{,}0\text{ ha}}$.
  \item \textbf{(E)}, Incorreta. A execução de um PRAD envolve implantação, manutenção, controle de invasoras, replantio de falhas e monitoramento contínuo até a comprovação do estabelecimento ecológico e estabilização física exigidos pelo órgão ambiental.
  \item \textbf{(A)}, Combina corretamente as bases abertas oficiais (TOPODATA/SRTM, MapBiomas, SICAR/CAR) com o sistema plano métrico SIRGAS 2000/UTM e os elementos cartográficos obrigatórios.
  \item \textbf{(C)}, $F_{\text{norm}} = (4 - 1)/4 = 0{,}75$; $(1 - IQS) = 1 - 0{,}52 = 0{,}48$. Logo: $IPI = 0{,}35(0{,}75) + 0{,}35(0{,}48) + 0{,}30(0{,}70) = 0{,}2625 + 0{,}1680 + 0{,}2100 = \mathbf{0{,}6405}$ (Prioridade Alta, pois $0{,}60 \le 0{,}6405 < 0{,}80$).
\end{enumerate}
\end{boxGabarito}

\begin{boxGabarito}[Espelho de Correção: Parte II (Questões Dissertativas 9 e 10)]
\textbf{Questão 9 (2,0 pontos):}
\begin{itemize}[leftmargin=1.4em, itemsep=2.5pt]
  \item \textbf{Item (a), Enquadramento legal (1,0 pt):} informar que, para rio de $15\text{ m}$ de largura (faixa de $10\text{ m}$ a $50\text{ m}$), a Lei nº 12.651/2012 (Art. 4º, I, ``b'') exige \textbf{$50\text{ m}$} de APP em cada margem. Explicar que a Lei nº 6.938/1981 (Art. 14, § 1º) adota a \textbf{responsabilidade civil objetiva} (teoria do risco integral), bastando a comprovação do dano e do nexo causal para obrigar o poluidor/empreendedor à reparação integral, independentemente de culpa ou dolo.
  \item \textbf{Item (b), Bioengenharia e PRAD (1,0 pt):} propor controle prévio do escoamento a montante, retaludamento suave ou proteção da base da margem (ex.: feixes vivos, estacas vivas e biomanta biodegradável no talude) associado ao plantio de espécies nativas ripárias adaptadas e isolamento contra pisoteio. Citar ao menos dois indicadores de monitoramento: taxa de sobrevivência de mudas ($\ge 80\%$), cobertura vegetal do solo ($\%$), variação de pinos de erosão ($\Delta h$) ou estabilidade de perfil transversal.
\end{itemize}

\textbf{Questão 10 (2,0 pontos):}
\begin{itemize}[leftmargin=1.4em, itemsep=2.5pt]
  \item \textbf{Item (a), Cálculo e classificação do $IPI$ (1,0 pt):}
  \textbf{Setor 1:} $F_{\text{norm}} = (5 - 1)/4 = 1{,}00$; $(1 - IQS) = 1 - 0{,}32 = 0{,}68$.
  $IPI_1 = 0{,}35(1{,}00) + 0{,}35(0{,}68) + 0{,}30(0{,}85) = 0{,}350 + 0{,}238 + 0{,}255 = \mathbf{0{,}843}$ (\textbf{Prioridade Crítica}, pois $IPI_1 \ge 0{,}80$). \\
  \textbf{Setor 2:} $F_{\text{norm}} = (2 - 1)/4 = 0{,}25$; $(1 - IQS) = 1 - 0{,}76 = 0{,}24$.
  $IPI_2 = 0{,}35(0{,}25) + 0{,}35(0{,}24) + 0{,}30(0{,}30) = 0{,}0875 + 0{,}0840 + 0{,}0900 = \mathbf{0{,}2615} \approx \mathbf{0{,}262}$ (\textbf{Prioridade Baixa}, pois $IPI_2 < 0{,}35$).
  \item \textbf{Item (b), Cartografia e proposta para o Setor 1 (1,0 pt):} citar ao menos três dos seis elementos cartográficos obrigatórios (título informativo, seta Norte, escala gráfica/numérica, grade de coordenadas com datum SIRGAS 2000, legenda e quadro de fontes/autoria) e recomendar intervenção imediata integrada para o Setor 1 (desvio/dissipação de escoamento a montante, paliçadas permeáveis ou biomanta, descompactação/adubação orgânica e revegetação com monitoramento).
\end{itemize}
\end{boxGabarito}

\section*{Tabela de Diagnóstico de Revisão do Módulo 5}

Registre seu desempenho no simulado para direcionar a revisão final antes da entrega do Trabalho Prático (Semana 16) e do Simulado Geral Final:

\begin{center}
\renewcommand{\arraystretch}{1.35}
\small
\begin{tabularx}{\textwidth}{@{}p{2.2cm} p{4.2cm} c c >{\raggedright\arraybackslash}X@{}}
  \toprule
  \rowcolor{uefsNavy}
  \textcolor{white}{\bfseries Capítulo} & \textcolor{white}{\bfseries Eixo Avaliado} & \textcolor{white}{\bfseries Questões} & \textcolor{white}{\bfseries Acertos} & \textcolor{white}{\bfseries Ação Recomendada se Abaixo de 70\%} \\
  \midrule
  \textbf{Cap. 13} & Bioengenharia, Mohr--Coulomb e Paliçadas & Q1, Q2, Q3, Q9b & \quad / 4 & Refazer exemplos de dimensionamento de paliçadas e revisar efeitos hidrológicos e mecânicos das raízes. \\
  \rowcolor{bgGray}
  \textbf{Cap. 14} & PNMA, Código Florestal, Licenciamento e PRAD & Q4, Q5, Q6, Q9a & \quad / 4 & Revisar faixas de APP (Art. 4º), Reserva Legal (Art. 12), LP/LI/LO e etapas do PRAD. \\
  \textbf{Cap. 15} & Diagnóstico Aplicado, Cartografia e $IPI$ & Q7, Q8, Q10 & \quad / 3 & Refazer Exemplos 15.1 a 15.5 e auditar seu Trabalho Prático Final pela rubrica de $5{,}0$ pontos. \\
  \bottomrule
\end{tabularx}
\end{center}
